%==============================================================
%  lecons/arithmetique/decomposition_premiers.tex
%  Théorème fondamental de l'arithmétique
%==============================================================

\lecontitle{Théorème fondamental de l'arithmétique}{Arithmétique — Nombres premiers et décomposition}

%=============================================================
%  § 1 — RAPPELS ET DÉFINITIONS
%=============================================================
\secnum{navyblue}{1}{Définitions et outils préliminaires}

\begin{tcolorbox}[thmbox]
  \textbf{Définitions.}\enspace%
  \index{nombre premier}\index{divisibilité}
  Un entier $p \geq 2$ est \emph{premier} s'il n'admet pas d'autre diviseur
  positif que $1$ et $p$.
  Un entier $n \geq 2$ non premier est dit \emph{composé}.

  \medskip
  \textbf{Lemme de Bézout.}\enspace\index{Bézout (identité de)}
  \textit{Pour tous $a, b \in \mathbb{Z}$, il existe $u, v \in \mathbb{Z}$ tels que}
  $au + bv = \pgcd(a,b)$.
  \textit{En particulier, $\pgcd(a,b) = 1 \Rightarrow \exists\,u,v : au+bv=1$.}

  \medskip
  \textbf{Lemme d'Euclide.}\enspace\index{Euclide (lemme d')}
  \textit{Soit $p$ premier. Si $p \mid ab$, alors $p \mid a$ ou $p \mid b$.}
\end{tcolorbox}

\rubriq{Preuve du lemme d'Euclide}

\begin{mdframed}[style=proofbar]
  Supposons $p \mid ab$ et $p \nmid a$. Comme $p$ est premier, $\pgcd(p,a) \in \{1,p\}$,
  et $p \nmid a$ impose $\pgcd(p,a)=1$. Par Bézout, il existe $u,v$ tels que
  $pu + av = 1$. En multipliant par $b$ :
  \[
    pub + avb = b.
  \]
  $p \mid pub$ (car $p \mid p$) et $p \mid avb$ (car $p \mid ab$), donc $p \mid b$.
  $\blacksquare$
\end{mdframed}

\rubriq{Généralisation}
Par récurrence : si $p$ est premier et $p \mid a_1 a_2 \cdots a_k$,
alors $p \mid a_i$ pour au moins un $i$.

\decsep{}

%=============================================================
%  § 2 — THÉORÈME FONDAMENTAL
%=============================================================
\secnum{navyblue}{2}{Théorème fondamental de l'arithmétique}

\begin{tcolorbox}[thmbox]
  \textbf{Théorème (décomposition en facteurs premiers).}\enspace%
  \index{décomposition en facteurs premiers}
  \textit{Tout entier $n \geq 2$ s'écrit de façon \textbf{unique} (à l'ordre des
  facteurs près) comme produit de nombres premiers :}
  \[
    n = p_1^{\alpha_1}\, p_2^{\alpha_2} \cdots p_k^{\alpha_k},
    \qquad p_1 < p_2 < \cdots < p_k \text{ premiers},\quad \alpha_i \geq 1.
  \]
\end{tcolorbox}

\rubriq{Preuve — Existence}

\begin{mdframed}[style=proofbar]
  Par récurrence forte sur $n$.
  \textit{Base :} $n=2$ est premier. \\
  \textit{Hérédité :} Soit $n \geq 3$. Si $n$ est premier, c'est sa propre décomposition.
  Sinon, $n = ab$ avec $2 \leq a,b < n$. Par hypothèse de récurrence, $a$ et $b$
  admettent chacun une décomposition en facteurs premiers. En les concaténant,
  on obtient une décomposition de $n$.
  $\blacksquare$
\end{mdframed}

\rubriq{Preuve — Unicité}

\begin{ideepreuve}
  C'est l'unicité, non l'existence, qui fait la profondeur du théorème : elle repose
  tout entière sur le lemme d'Euclide. L'idée est d'éplucher les deux décompositions
  un facteur à la fois : le lemme d'Euclide garantit qu'un premier $p_1$ du premier
  membre doit apparaître à l'identique dans le second, on le simplifie, et l'on
  recommence sur un entier plus petit.
\end{ideepreuve}

\begin{mdframed}[style=proofbar]
  Supposons $n = p_1^{\alpha_1}\cdots p_k^{\alpha_k} = q_1^{\beta_1}\cdots q_l^{\beta_l}$
  avec $p_i, q_j$ premiers. On montre par récurrence sur
  $N = \alpha_1 + \cdots + \alpha_k$, nombre de facteurs de la \emph{première}
  décomposition comptés avec multiplicité, que les deux décompositions coïncident.
  (Une fois l'unicité acquise, ce nombre ne dépend plus de la décomposition choisie :
  on le note alors $\Omega(n)$.)

  \textit{Base :} $N=1$ signifie que $n = p_1$ est premier. Toute autre
  décomposition $n = q_1^{\beta_1}\cdots q_l^{\beta_l}$ a au moins un facteur ;
  si elle comportait au moins deux facteurs premiers (comptés avec multiplicité),
  alors $n$ s'écrirait comme produit de $\ge 2$ premiers, donc admettrait un diviseur
  premier strictement compris entre $1$ et $n$, ce qui contredit la primalité de $n$.
  Ainsi $l=1$, $\beta_1=1$ et $q_1 = n = p_1$ : l'unicité est acquise.

  \textit{Hérédité :} $p_1$ divise $q_1^{\beta_1}\cdots q_l^{\beta_l}$, donc par le
  lemme d'Euclide généralisé, $p_1$ divise l'un des $q_j$ ; comme $q_j$ est premier
  et $p_1 \geq 2$, on a $p_1 = q_j$.
  Quitte à réordonner, $p_1 = q_1$. On divise par $p_1$ :
  $n/p_1 = p_1^{\alpha_1-1}\cdots p_k^{\alpha_k} = q_1^{\beta_1-1}\cdots q_l^{\beta_l}$.
  La première décomposition de $n/p_1$ compte $N-1$ facteurs ; par hypothèse de
  récurrence, les deux factorisations de
  $n/p_1$ coïncident, donc $k=l$, $p_i=q_i$ et $\alpha_i=\beta_i$ pour tout $i$.
  $\blacksquare$
\end{mdframed}

\decsep{}

%=============================================================
%  § 3 — CONSÉQUENCES ET APPLICATIONS
%=============================================================
\secnum{navyblue}{3}{Conséquences}

\begin{tcolorbox}[thmbox]
  Soient $a = \prod p_i^{\alpha_i}$ et $b = \prod p_i^{\beta_i}$
  (en complétant par des exposants nuls).
  \begin{enumerate}
    \item $\pgcd(a,b) = \prod p_i^{\min(\alpha_i,\beta_i)}$,
          $\;\mathrm{ppcm}(a,b) = \prod p_i^{\max(\alpha_i,\beta_i)}$.
    \item $\pgcd(a,b) \cdot \mathrm{ppcm}(a,b) = ab$.
    \item $a \mid b \iff \alpha_i \leq \beta_i$ pour tout $i$.
    \item \textbf{Infinité des nombres premiers (Euclide).}\enspace%
    \index{nombres premiers!infinité}
    L'ensemble des nombres premiers est infini.
  \end{enumerate}
\end{tcolorbox}

\rubriq{Preuve de l'infinité des premiers}

\begin{mdframed}[style=proofbar]
  Supposons par l'absurde qu'il n'existe qu'un nombre fini de premiers $p_1,\ldots,p_k$.
  Posons $N = p_1 p_2 \cdots p_k + 1$. Comme $N \geq 2$, le théorème lui fournit un
  facteur premier $q$. Puisque $p_1,\ldots,p_k$ sont \emph{tous} les premiers, $q$
  figure dans cette liste, d'où $q \mid p_1\cdots p_k$. Mais $q \mid N$ ; donc $q$
  divise la différence $N - p_1\cdots p_k = 1$, ce qui force $q \leq 1$ :
  contradiction.
  $\blacksquare$
\end{mdframed}

\decsep{}

%=============================================================
%  § 4 — APPLICATIONS, PIÈGES, EXERCICES
%=============================================================
\secnum{exgreen}{4}{Applications, pièges et exercices}

\rubriq{Applications}

\begin{mdframed}[style=exbar]
  \textbf{1. Irrationalité de $\sqrt{2}$.}\enspace%
  \index{irrationalité!de $\sqrt{2}$}
  Si $\sqrt{2} = p/q$ avec $\pgcd(p,q)=1$, alors $p^2 = 2q^2$.
  Dans la décomposition de $p^2$, la puissance de $2$ est paire ;
  dans celle de $2q^2$, elle est impaire. Contradiction.

  \smallskip\noindent
  \textbf{2. Nombre de diviseurs.}\enspace%
  Si $n = \prod p_i^{\alpha_i}$, alors $n$ admet $\prod(\alpha_i+1)$ diviseurs positifs.

  \smallskip\noindent
  \textbf{3. Valuation $p$-adique.}\enspace%
  \index{valuation $p$-adique}
  La valuation $v_p(n)$ = exposant de $p$ dans la décomposition de $n$ vérifie
  $v_p(ab) = v_p(a)+v_p(b)$ et $v_p(a+b) \geq \min(v_p(a),v_p(b))$.
\end{mdframed}

\vspace{6pt}
\rubriq{Pièges}

\begin{mdframed}[style=cexbar]
  \textbf{L'unicité dépend du lemme d'Euclide, pas de la seule existence.}\enspace%
  Dans $\mathbb{Z}[\sqrt{-5}]$, on a $6 = 2\times 3 = (1+\sqrt{-5})(1-\sqrt{-5})$,
  deux factorisations en « irréductibles » différentes. L'unicité est une propriété
  spécifique à $\mathbb{Z}$ (anneau factoriel, et même principal).

  \smallskip\noindent
  \textbf{$\pgcd(a,b)=1$ ne signifie pas que $a$ ou $b$ est premier.}\enspace%
  $\pgcd(4,9)=1$ mais $4$ et $9$ sont composés.
\end{mdframed}

\vspace{6pt}
\exolabel

\begin{mdframed}[style=exobar]
  \begin{enumerate}
    \item Montrer que $\sqrt{p}$ est irrationnel pour tout premier $p$.
          Généraliser : $\sqrt[k]{n}$ est irrationnel si $n$ n'est pas une
          puissance $k$-ième parfaite.

    \item Calculer $\pgcd(2^{100}-1, 2^{75}-1)$.
          (\textit{Indication :} montrer que $\pgcd(2^a-1, 2^b-1) = 2^{\pgcd(a,b)}-1$.)

    \item Montrer que pour tout $n \geq 1$, le produit de $n$ entiers consécutifs
          est divisible par $n!$. (\textit{Indication :} coefficients binomiaux.)

    \item (Théorème de Wilson) Montrer qu'un entier $p \geq 2$ est premier si et
          seulement si $(p-1)! \equiv -1 \pmod{p}$.

    \item Montrer que si $2^n - 1$ est premier, alors $n$ est premier
          (les premiers de cette forme sont dits de Mersenne\index{nombre de Mersenne}).
  \end{enumerate}
\end{mdframed}

\leconfooter{Euclide, \textit{Éléments} (vers $-300$) — C.\ F.\ Gauss, \textit{Disquisitiones} (1801)}
